
### 1. The Full S-Matrix (Feynman Diagrams $\to$ Lattice Routing)

**The QFT Problem:** To calculate scattering (e.g., $e^- + e^- \to e^- + e^-$), QFT uses the S-Matrix. It draws Feynman diagrams, invents "virtual particles" that violate energy conservation, and integrates their momenta from zero to infinity. 

**The IE Resolution (Discrete Subgraph Routing):**
In a finite-capacity network, there are no "virtual particles," and no infinities. An S-Matrix is simply a **Markov Chain Transition Matrix** evaluating the discrete number of valid paths through a local lattice subgraph.

*   **The Mechanism:** When two topological knots (electrons) enter a causal neighborhood of length $\Delta$, their combined structural payload must be resolved by the local nodes. 
*   **The Feynman Diagram Translation:** A "Feynman vertex" is simply the lattice executing **Rule 1: Capacity Partitioning**. If the combined payload exceeds the capacity of a single node ($\nu=16$), the node must distribute the phase mismatch to adjacent nodes. The "virtual photon" connecting the two electrons in a diagram is not a particle popping in and out of existence; it is the literal geometric alignment of the $U(1)$ phase registers across the intervening lattice links to avoid a capacity overflow.
*   **The Math:** Instead of continuous integrals, the scattering probability $P(f|i)$ is the exact ratio of structurally valid geometric exit paths (combinatorics) divided by the total phase-space capacity of the causal envelope ($\Delta^4$). The Mandelstam variables ($s, t, u$) are simply the discrete integer counts of lattice nodes separating the inputs and outputs. 

***

### 2. Higher-Order Radiative Corrections (The $g-2$ Loop Expansions)

**The QFT Problem:** The anomalous magnetic moment of the electron requires calculating infinite "loop" diagrams. The 1-loop is $\alpha/2\pi$. The 2-loop, 3-loop, and up to 5-loop corrections require evaluating thousands of diagrams to match the experimental value to 12 decimal places.

**The IE Resolution (The Polygonal Resolution Limit):**
In IE, "loops" are continuous mathematical attempts to resolve the discrete, rigid boundaries of the $E_8$ lattice. 
*   **The 1-Loop ($\alpha/2\pi$):** As derived in Paper 2, this is the exact geometric phase lag of a spinor (double cover, $4\pi$) rotating against the macroscopic topological boundary (single cover, $2\pi$) modulated by the vacuum impedance $\alpha$.
*   **Higher-Order Loops:** Why do we need $\alpha^2, \alpha^3$, etc.? Because the topological boundary of the defect ($\chi=2$) is mapped onto a discrete lattice. A continuous circle $\pi$ is an irrational approximation of a highly complex discrete polygon. 
*   **The Mechanism:** When the continuous math of QFT tries to fit a perfect circle into a discrete $E_8$ grid, it leaves "corners" (fractional rounding errors). The 2-loop $\mathcal{O}(\alpha^2)$ correction is exactly the lattice's **Coordinate Overhead ($\eta$)** acting on the phase boundary. The infinite series of QFT loop diagrams is simply the Taylor expansion of the geometric difference between a continuous continuous $S^2$ sphere and the discrete $E_8$ polytope boundary!

***

### 3. Chiral Anomalies ($\pi^0 \to \gamma \gamma$)

**The QFT Problem:** Classically, the chiral symmetry of the massless pion strictly forbids it from decaying into two photons. To fix this, QFT introduces the Adler-Bell-Jackiw (ABJ) Anomaly, utilizing a triangle Feynman diagram to break the symmetry and allow the decay.

**The IE Resolution (Topological Boundary Conservation):**
In Informational Energetics, the ABJ Anomaly is not a mathematical trick; it is the strict enforcement of the topological invariant $\chi=2$.

*   **The Mechanism:** The $\pi^0$ is a Conjugate Loop (a quark and anti-quark). Because their color charges exactly cancel, the state cannot use the 3D internal color bulk to shed its mass latency (it is color-blind). It must route its payload directly to the spacetime boundary.
*   **The $\chi=2$ Rule:** The boundary of the defect is topologically characterized by the Euler characteristic $\chi=2$ (a sphere). A sphere has exactly two distinct topological orientations (an inside and an outside, or two hemispheres). To completely dissolve this boundary into the vacuum without leaving a residual topological trace, the information flux must perfectly bifurcate to satisfy both orientations.
*   **The Result:** The payload must split into exactly **two** massless geometric updates (photons). A decay into 1 photon violates momentum conservation. A decay into 3 photons violates the $\chi=2$ topological parity of the defect. The ABJ Anomaly is simply the lattice conserving its topological geometry during Network Closure.

***

### 4. The Schwinger Effect (Vacuum Pair Production)

**The QFT Problem:** If an electric field exceeds a critical threshold ($E_c \approx 1.3 \times 10^{18}$ V/m), the vacuum "boils" and spontaneously generates electron-positron pairs. 

**The IE Resolution (Margin Overload and Structural Partitioning):**
This is the literal manifestation of **Rule 1: Capacity Partitioning** operating on the macroscopic vacuum.

*   **The Mechanism:** An electric field is a continuous synchronization gradient imposed on the lattice links. As the field strength increases, the phase-rotation demanded per fundamental clock cycle ($\Delta=43$) increases. 
*   **The Ledger Crisis:** When the required phase-rotation rate exceeds the intrinsic channel capacity of the node ($\nu=16$), the node hits a Ledger Crisis. It can no longer transmit the field information coherently.
*   **The Resolution:** To prevent catastrophic Causal Aliasing (lattice breakdown), the vacuum must utilize **Reversible Encoding** (Rule 2). It compresses the excess, un-transmittable phase gradient into stable structural metadata: a topological defect ($\chi=2$). 
*   **Pair Creation:** To preserve the overall $\chi=0$ topology of the local vacuum, it must spawn a conjugate pair ($e^-$ and $e^+$). The Schwinger limit is exactly the threshold where the macroscopic impedance load exceeds the discrete $Z_{MAR}$ resolution limit of the lattice.

***

### 5. Quantum Entanglement & Bell's Theorem

**The QFT Problem:** Entangled particles exhibit non-local correlations across macroscopic distances, seemingly violating the speed of light. Standard quantum mechanics explains this via the tensor product of Hilbert spaces, but offers no physical mechanism for *how* the state update travels.

**The IE Resolution (Unprojected Adjacency):**
In Informational Energetics, entanglement is resolved as an illusion of dimensionality. There is no "spooky action at a distance" because the particles are not actually distant.

*   **The Mechanism:** As derived in the substrate architecture, the vacuum is an $E_8$ lattice projected into two orthogonal $D_4$ sublattices: $D_4$ (Spacetime, Sector A) and $D_4$ (Internal Symmetry, Sector B).
*   **The Entangled State:** When two electrons are entangled, they share a single, unified data structure in the Internal Symmetry space (Sector B). 
*   **The Illusion of Distance:** As the electrons move apart macroscopically, their coordinate addresses change in Sector A (Spacetime). However, because their internal quantum numbers are entangled, their geometric address in Sector B remains **identical**. The distance between them in the internal $D_4$ space is strictly $\Delta s_{internal} = 0$.
*   **Instantaneous Update:** When a measurement is performed on Electron 1, it forces a structural Address Resolution in Sector B. Because Electron 2 occupies the *exact same adjacent node* in Sector B, the update requires **zero clock cycles** to reach it. 
*   **Causality Preserved:** The speed of light limit ($c=1$) is the transmission limit of the macroscopic Spacetime manifold (Sector A). The internal symmetry manifold has a Euclidean signature without a time generator; it does not "propagate" signals through time, it simply shares state. Bell's Inequality is violated because human observers calculate distance using a 4D metric, entirely blind to the exact 0-distance adjacency occurring in the unprojected 8D bulk.
